% =====================================================================
\chapter{Multi-Layer Perceptron}
% =====================================================================

\labinfo{08\_MLP}{mlp.py}{UCI Iris: 150 flowers, 4 standardised measurements}{3-class classification
(setosa / versicolor / virginica)}{Test accuracy 96.7\%, macro F1 0.966 (4 hidden units)}

\section{Dataset}
The \textbf{Iris} dataset (UCI) contains 150 flowers, 50 from each of three species. Each flower has
four numeric measurements: sepal length, sepal width, petal length and petal width. I standardised the
four inputs and split the data into 98 training, 22 validation and 30 test flowers.

\section{Code and experiments}
\begin{itemize}
\item \code{sigmoid}, \code{softmax} (which subtracts the row maximum before exponentiating, to avoid
      overflow) and \code{one\_hot} are small helpers.
\item The \code{MLP} class holds $W_1$ (64×H), $\vect b_1$, $W_2$ (H×10) and $\vect b_2$.
      \code{forward} returns $\vect h$ and $\vect p$. \code{backward} implements the four
      backpropagation formulas above and applies the gradient step. \code{fit} runs mini-batch training
      (batch 32, learning rate 0.5, 100 epochs) and records the learning curves.
\item \code{confusion\_matrix} uses \code{np.add.at} to count (true, predicted) pairs, and
      \code{per\_class\_scores} turns the matrix into precision, recall and F1 per digit.
\end{itemize}
\noindent Experiments: (1)~hidden sizes 4, 8, 16, 32, 64 and 128, compared on the validation set;
(2)~the chosen network trained once more and evaluated on the test file.

\section{Results}

\begin{table}[H]
\centering\small
\caption{Choosing the size of the hidden layer (100 epochs each).}
\label{tab:mlp-hidden}
\begin{tabular}{@{}l c c c c c c@{}}
\toprule
Hidden units & 4 & 8 & 16 & 32 & 64 & 128\\
Weights      & 310 & 610 & 1,210 & 2,410 & 4,810 & 9,610\\
\midrule
Training accuracy   & 0.959 & 0.959 & 0.959 & 0.888 & 0.806 & 0.684\\
Validation accuracy & \textbf{0.955} & 0.864 & 0.818 & 0.727 & 0.591 & 0.591\\
\bottomrule
\end{tabular}
\end{table}

Four hidden units gave the best validation result on this small Iris split, so the final network uses
4 inputs, 4 hidden units and 3 outputs (310 parameters). Larger hidden layers overfit the small training
sample and performed worse on validation.

\begin{table}[H]
\centering\small
\caption{Test-set results (30 Iris flowers).}
\label{tab:mlp-test}
\begin{tabular}{@{}l c c c c c c c c c c@{}}
\toprule
\multicolumn{11}{@{}l}{Accuracy \textbf{0.9667} (1 mistake) \quad Macro precision 0.9722 \quad Macro recall 0.9630 \quad Macro F1 0.9659}\\
\midrule
Species & Setosa & Versicolor & Virginica \\
F1 & 1.000 & 0.941 & 0.957\\
\bottomrule
\end{tabular}
\end{table}

The network recognises 96.7\% of the test flowers. The single error is between the two species whose
measurements overlap most, while setosa is separated cleanly by the petal measurements.

\begin{figure}[H]
\centering
\includegraphics[width=0.48\linewidth]{\figpath{08\_MLP}{learning_curves.pdf}}\hfill
\includegraphics[width=0.48\linewidth]{\figpath{08\_MLP}{confusion_matrix.pdf}}
\caption{MLP learning curves and Iris test confusion matrix.}
\end{figure}

\section{Source code}
\lstinputlisting[caption={\texttt{08\_MLP/mlp.py}}]{08_MLP/mlp.py}
\clearpage
